\section{Related Work}
\label{sec:related-work}

\paragraph{Self-Consistency and Extensions.}
Self-Consistency (SC) selects the most frequent answer from multiple sampled chain-of-thought traces \citep{wang2023sc} and remains a strong baseline for multi-path reasoning. Universal Self-Consistency (USC) extends SC to free-form generation by having the LLM itself select the most consistent response \citep{chen2024usc}. Confidence-weighted variants assign trace-level confidence scores to improve vote quality \citep{taubenfeld2025cisc}, while DiVeRSe combines diverse prompting strategies with step-aware verification \citep{li2023diverse}. Recent work has identified scaling plateaus: increasing the number of sampled traces yields diminishing returns beyond a saturation point \citep{loo2025reevaluating}. These findings motivate the search for aggregation mechanisms that extract richer signal from reasoning traces than simple vote counting. SICA represents one such attempt: it replaces plurality voting with symbolic constraint satisfaction over self-extracted FOL formulas. Our negative results show that the bottleneck lies not in the aggregation mechanism but in the quality of self-extracted constraints, a failure mode that vote-based methods sidestep by not requiring formalization at all.

\paragraph{Neuro-Symbolic Reasoning.}
A growing line of work combines LLMs with symbolic solvers to improve logical reasoning. Logic-LM translates natural language problems into symbolic formulations and delegates solving to an external prover \citep{pan2023logiclm}. LINC pairs language models with first-order logic provers in a modular architecture \citep{olausson2023linc}. More recent systems push this integration further: DyLoS adaptively selects among multiple solver paradigms \citep{xu2026dylos}, Aristotle decomposes problems through a logic-complete search framework \citep{li2025aristotle}, FoVer trains a dedicated verifier on FOL annotations \citep{wang2025fover}, and LogicAgent applies multi-perspective semiotic grounding \citep{luo2025logicagent}. A critical distinction separates these methods from SICA. All of the above rely on trained verifiers, external provers, or multi-paradigm solver selection, components that introduce knowledge beyond the reasoning model itself. SICA instead asks whether the same model that produced the reasoning traces can also extract useful formal constraints from them. Our oracle experiment confirms that sound FOL constraints do improve aggregation (+8.82\,pp); the failure is specific to same-model self-extraction, not to symbolic reasoning as an approach.

\paragraph{Test-Time Compute and Reasoning Models.}
Scaling test-time computation has emerged as an alternative to parameter scaling. \citet{snell2025scaling} show that optimally allocated test-time compute can outperform larger models and establish a framework for analyzing the compute-accuracy tradeoff at inference. Reasoning-specialized models such as DeepSeek-R1 \citep{deepseek2025r1} internalize extended chain-of-thought through reinforcement learning and produce longer but more coherent traces. Diversity-aware training methods explicitly optimize for trace diversity during test-time scaling \citep{diverseaware2025}. These developments interact with our findings through the $C_0$ diagnostic. Reasoning models produce highly consistent outputs across samples: $C_0$ drops below 5\% for DeepSeek-R1 and QwQ on our benchmarks. When nearly all traces converge to the same answer, any aggregation method (symbolic or otherwise) reduces to the unanimous output. The $C_0$ metric thus provides a zero-cost signal for practitioners to determine whether post-hoc aggregation can offer any benefit for a given model.

\paragraph{Verification and Reranking.}
Trained verifiers offer an alternative path to improving multi-path reasoning. Outcome-level reward models (ORMs) score complete solutions and select the highest-ranked candidate \citep{cobbe2021gsm8k}. Process reward models (PRMs) assign step-level scores and enable finer-grained selection \citep{lightman2023prm}. Math-Shepherd automates process supervision annotation and removes the need for human labels \citep{wang2024mathshepherd}. These approaches share SICA's goal of selecting better answers from a candidate pool but differ in a fundamental respect: they use a separately trained scoring model. SICA's constraint extraction requires no training data and no additional model; it relies entirely on the reasoning model's own ability to formalize its logic. Our results suggest that this self-reliance is precisely the problem. Without an external signal to ground constraint quality, self-extracted constraints encode the same biases present in the original traces and fail to provide corrective information.

\paragraph{Constraint Satisfaction in NLP.}
Constraint-based formulations have gained traction for evaluating and improving LLM reasoning. SATBench benchmarks LLMs on Boolean satisfiability problems and reveals systematic weaknesses in logical structure handling \citep{satbench2025}. Work on integrating LLMs with constraint propagation solvers explores hybrid architectures where the LLM generates candidate assignments and a CP solver enforces hard constraints \citep{llmmeets_cp2025}. SICA's MAX-SAT formulation sits within this broader trend: it uses Z3 to find the maximum-weight satisfiable subset of self-extracted constraints. The key difference is that SICA's constraints originate from the model's own reasoning traces rather than from problem specifications or external knowledge bases. Our finding that these self-extracted constraints lack discriminative power complements the SATBench results: LLMs struggle both to solve formal constraints directly and to produce useful constraints from their own reasoning.

\paragraph{Debate and Multi-Agent Aggregation.}
Multi-agent debate refines answers through iterative natural language discussion among LLM instances \citep{liang2023encouraging,du2023improving}. In debate, models exchange arguments in plain text and update their positions based on peer responses. SICA takes a structurally different approach: instead of natural language refinement, it extracts formal constraints and resolves disagreements through satisfiability optimization. Both families attempt to move beyond simple voting, but debate preserves natural language flexibility while SICA commits to formal structure. \citet{huang2024selfcorrect} show that LLMs struggle to self-correct reasoning without external feedback, a finding that parallels our capability trap: the same model that produced flawed reasoning cannot reliably formalize the correction. Across all six lines of work reviewed here, the consistent pattern is that improving on majority voting requires information beyond what the reasoning model can extract from its own outputs, whether from trained verifiers, external solvers, human annotations, or cross-model interaction.
